\documentclass[11pt]{article}
\usepackage[margin=1in]{geometry}
\usepackage{amsmath,amssymb,amsthm,mathtools}
\usepackage[hidelinks]{hyperref}

\newtheorem{theorem}{Theorem}[section]
\newtheorem{proposition}[theorem]{Proposition}
\newtheorem{remark}[theorem]{Remark}
\newcommand{\splus}{s^+}
\newcommand{\sminus}{s^-}
\newcommand{\sig}{\sigma}
\newcommand{\one}{\mathbf 1}

\title{Positive Square Energy of Every Nonunicyclic Cactus}
\author{AI-generated manuscript; no human author is claimed}
\date{30 July 2026}

\begin{document}
\maketitle

\begin{abstract}
We prove that every finite simple connected cactus $G$ of cyclomatic rank
$k\geq2$ and order $n$ satisfies $s^+(G)>n$.  The exact cactus doubly
nonnegative bound gives
\[
 s^+(G)-n\geq k-1-\sum_i\epsilon_{\ell_i}.
\]
A rank-independent deficiency calculation shows that its non-strict frontier
is exactly $T^{k-1}Q$ and $T^{k-2}PP$, where $T=C_3$, $P=C_5$, and $Q$ is an
arbitrary cycle.  Existing hostile and nonhostile all-rank one-cycle theorems
close the first family.  The all-rank two-pentagon theorem closes the second
for $k\geq3$, while at $k=2$ the pure two-pentagon family follows from the
separate shared-cut and positive-connector theorems.  Rank one is explicitly
excluded: $C_4$ has $s^+(C_4)=4=|V(C_4)|$.  A constant-size, fail-closed exact
verifier freezes the uniform frontier calculation.
\end{abstract}

\section{Statement and the rank-one boundary}

All graphs in this paper are finite, simple, and undirected.  A cactus is
connected and has only bridge or cycle blocks.  Its cyclomatic rank is the
number of cyclic blocks.  For the adjacency eigenvalues
$\lambda_1,\ldots,\lambda_n$, put
\[
 \splus(G)=\sum_{\lambda_i>0}\lambda_i^2,
 \qquad
 \sminus(G)=\sum_{\lambda_i<0}\lambda_i^2,
 \qquad
 \sig(G)=\splus(G)-|V(G)|.
\]

\begin{theorem}\label{thm:main}
Every finite simple connected cactus $G$ of cyclomatic rank $k\geq2$ satisfies
\[
 \boxed{\splus(G)>|V(G)|}.
\]
\end{theorem}

The lower bound $k\geq2$ is essential for strictness.  The adjacency spectrum
of $C_4$ is $2,0,0,-2$, and therefore
\[
 \splus(C_4)=4=|V(C_4)|.
\]
Thus Theorem~\ref{thm:main} neither includes nor can be extended verbatim to
all unicyclic cacti.

\section{The rank-uniform sharp-DNN frontier}

For an edge-containing graph define
\[
 \kappa(G)=\sup_{\substack{M\succeq0,\ M\geq0\\\one^TM\one>0}}
 \frac{4\bigl(\sum_{uv\in E(G)}\sqrt{M_{uv}}\bigr)^2}
      {\one^TM\one}.
\]
The exact cactus DNN theorem~\cite{DNN,LTZ} says that if a cactus has $b$
bridge blocks and cycle blocks $C_{\ell_1},\ldots,C_{\ell_k}$, then
\begin{equation}\label{eq:kappa}
 \sminus(G)\leq\kappa(G)
 =b+\sum_{i=1}^k\kappa(C_{\ell_i}),
 \qquad
 \kappa(C_\ell)=
 \begin{cases}
  \ell,&\ell\text{ even},\\
  \ell\sec^2\!\bigl(\pi/(2\ell)\bigr),&\ell\text{ odd}.
 \end{cases}
\end{equation}
Set
\[
 \epsilon_\ell=
 \begin{cases}
 0,&\ell\text{ even},\\
 \ell\tan^2\!\bigl(\pi/(2\ell)\bigr),&\ell\text{ odd}.
 \end{cases}
\]
If $n=|V(G)|$, block counting gives
$b+\sum_i\ell_i=n+k-1$, while
$\splus(G)+\sminus(G)=2|E(G)|=2(n+k-1)$.  Hence
\begin{equation}\label{eq:surplus}
 \boxed{\sig(G)\geq k-1-\sum_{i=1}^k\epsilon_{\ell_i}.}
\end{equation}

Write $T=C_3$, $P=C_5$, and
\[
 a=\epsilon_5=5-2\sqrt5.
\]
For real $x\geq3$, $x\tan^2(\pi/(2x))$ decreases strictly.  Indeed, with
$u=\pi/(2x)$, differentiation reduces the sign to
$2u-\sin u\cos u>0$.  Thus $\epsilon_3=1$ and every nontriangle contributes
at most $a$.  We use the exact comparisons
\begin{equation}\label{eq:exact}
 0<a<1,\qquad3a<2,\qquad2a>1,
 \qquad a+\epsilon_7<1.
\end{equation}
The last follows from $\cos(\pi/7)>2159/2401>7/8$, which gives
$\epsilon_7<7/15$, and from
$7/15<2\sqrt5-4=1-a$, certified by $(67/30)^2<5$.

\begin{proposition}[Exact uniform deficiency frontier]\label{prop:frontier}
For every integer $k\geq2$, the right side of~\eqref{eq:surplus} is not
strictly positive exactly on the two cycle-multiset families
\[
 \boxed{T^{k-1}Q},\qquad \boxed{T^{k-2}PP},
\]
where $Q=C_q$ is arbitrary and may equal $T$.
\end{proposition}

\begin{proof}
Let $h$ be the number of nontriangular cycles.  If $h\geq3$, monotonicity and
\eqref{eq:exact} give
\[
 \sum_i\epsilon_{\ell_i}\leq(k-h)+ha
 \leq(k-3)+3a<k-1.
\]
If $h=2$, a pair containing an even cycle contributes less than one.  If both
cycles are odd but are not both pentagons, their contribution is at most
$a+\epsilon_7<1$.  The sole remaining pair is $PP$, and it is genuinely
residual because $2a>1$.  This gives $T^{k-2}PP$.

If $h\leq1$, the multiset is $T^{k-1}Q$, allowing $Q=T$ when $h=0$.  Conversely
every such multiset has epsilon sum
$(k-1)+\epsilon_q\geq k-1$, so~\eqref{eq:surplus} is not strict.  The cases
are exhaustive.
\end{proof}

This proof has a fixed number of cases and inequalities independent of $k$.
The exact verifier~\cite{Frontier} represents each quantity as an affine
element $cK+d$ of $\mathbb Q(\sqrt5)[K]$.  It proves that the $K$ coefficient
cancels in every comparison and checks the exact sign of the remaining
constant.  It also proves that the $h\geq3$ bound decreases with $h$, so its
maximum is attained at $h=3$.  Substitutions at
$K=2,3,4,5,7,13,64$ semantically test every generated affine template, but no
finite collection of ranks is used to prove the arbitrary-rank claim.

\section{Closure of the two residual families}

\begin{proposition}[One distinguished cycle]\label{prop:q}
Every connected cactus with cycle multiset $T^{k-1}Q$, $k\geq2$, has positive
surplus.
\end{proposition}

\begin{proof}
There are $k-1\geq1$ triangles.  If $q\equiv1\pmod4$, then $q\geq5$ and the
all-rank triangle--hostile-cycle theorem applies~\cite{Hostile}.  If $q$ is
even or $q\equiv3\pmod4$, the all-rank nonhostile one-cycle theorem applies
\cite{Nonhostile}.  The latter explicitly includes $q=3$, hence also the
all-triangle member.  These alternatives exhaust every integer $q\geq3$.
Both inputs permit arbitrary shared cuts, bridge connectors, cactus incidence,
and finite tree attachments.
\end{proof}

\begin{proposition}[Two pentagons]\label{prop:pp}
Every connected cactus with cycle multiset $T^{k-2}PP$, $k\geq2$, has positive
surplus.
\end{proposition}

\begin{proof}
Suppose first that $k\geq3$.  Then $r=k-2\geq1$, and the all-rank theorem for
$r$ triangles and exactly two pentagons applies directly~\cite{TwoPent}.

It remains to respect the boundary $k=2$, where there are no triangles and the
all-rank $T^rPP$ theorem is outside its hypothesis.  If the two pentagons meet,
the cactus property forces their intersection to be one cut vertex; the
two-pentagon bouquet theorem gives
\[
 \splus(G)\geq n+1-\frac4{3\sqrt{13}}>n
\]
with arbitrary attached trees~\cite{Bouquet}.  If they do not meet, their
unique connecting path has positive length; the two-actual-lobes connector
theorem gives
\[
 \splus(G)>n+5-2\sqrt5>n
\]
for every connector length and arbitrary trees~\cite{Connector}.  The two
possibilities exhaust the pure $PP$ bicyclic family.
\end{proof}

\begin{proof}[Proof of Theorem~\ref{thm:main}]
Outside the two families of Proposition~\ref{prop:frontier},
equation~\eqref{eq:surplus} is strictly positive.  Proposition~\ref{prop:q}
closes $T^{k-1}Q$ for every $q\geq3$, and Proposition~\ref{prop:pp} closes
$T^{k-2}PP$, including its separate pure bicyclic boundary.  The frontier is
exhaustive.
\end{proof}

\section{Dependency map and reproducibility}

The logical dependency map is exact:
\begin{verbatim}
main theorem
|-- sharp DNN bound and uniform frontier
|   |-- sharp-cactus-dnn/paper.tex
|   `-- research/rank-uniform-cactus-dnn-frontier-verifier.py
|-- T^(k-1)Q
|   |-- all-rank-triangle-hostile-cacti/paper.tex
|   `-- research/all-rank-nonhostile-one-cycle-theorem-2026-07-30.md
`-- T^(k-2)PP
    |-- k>=3: all-rank-triangle-two-pentagon-cacti/paper.tex
    |-- k=2 shared cut: two-c5-bouquet-trees/paper.tex
    `-- k=2 connector: two-c5-all-connectors/paper.tex
\end{verbatim}
The earlier fixed-rank cactus manuscripts are corroborating specializations,
not premises.  A prose proof object with the same map is~\cite{ProofNote}.

From the repository root, reproduce the exact frontier and PDF with Python
3.10 or newer:
\begin{verbatim}
python3 research/rank-uniform-cactus-dnn-frontier-verifier.py
python3 -O research/rank-uniform-cactus-dnn-frontier-verifier.py
bash scripts/build-paper.sh all-nonunicyclic-cacti
\end{verbatim}
The two verifier outputs must be byte-identical.  The verifier uses exact
\texttt{Fraction} arithmetic in the affine ring $\mathbb Q(\sqrt5)[K]$ and
freezes a canonical digest of the generated symbolic expression records,
rather than prose or calls to a finite-rank frontier function.  Its certificate
SHA-256 digest is
\begin{center}
\texttt{1b549165c84ab20bdaebfec9329c37fb3e808ef435bda7a3586c8285c2d075ca},
\end{center}
and rejects seven hostile mutations: changing a $K$ coefficient, reversing the
radical sign in $a$, changing either the $h\geq3$ or $h=2$ boundary, or
deleting, reordering, or altering a survivor family.  It certifies only
Proposition~\ref{prop:frontier}; the DNN theorem and structural closures remain
the displayed mathematical dependencies.

\paragraph{Scoped nonclaim.}
Theorem~\ref{thm:main} concerns only finite simple connected cacti of
cyclomatic rank at least two.  It does not prove a statement for all graphs,
disconnected graphs, non-cactus block intersections, or unicyclic cacti.  In
particular, no universal all-graphs version of the positive-square-energy
conjecture is claimed.

\section*{AI disclosure}
This manuscript and its proof-note assembly were generated by an AI coding and
mathematical reasoning system from the cited repository proof objects.  No
human authorship, review, or verification is claimed.  The result should be
assessed from the displayed reduction and dependencies.  No external
publication was performed as part of this work.

\begin{thebibliography}{99}
\bibitem{LTZ} Y.~Liu, Q.~Tang, and S.~Zhang,
\emph{The positive and negative square-energy conjecture}, arXiv:2607.18031,
2026.

\bibitem{DNN} Companion manuscript, \emph{The Optimized Liu--Tang--Zhang
Constant for Cactus Graphs}, 2026; \path{sharp-cactus-dnn/paper.tex}.

\bibitem{Frontier} \emph{Rank-uniform cactus DNN frontier verifier}, 30 July
2026; \path{research/rank-uniform-cactus-dnn-frontier-verifier.py}.

\bibitem{ProofNote} \emph{Every nonunicyclic cactus has strict positive square
energy}, proof note, 30 July 2026;
\path{research/rank-uniform-cactus-theorem-proof-note-2026-07-30.md}.

\bibitem{Hostile} Companion manuscript, \emph{Cacti with Arbitrarily Many
Triangles and One Hostile Cycle}, 2026;
\path{all-rank-triangle-hostile-cacti/paper.tex}.

\bibitem{Nonhostile} \emph{All-rank nonhostile one-cycle cactus theorem},
30 July 2026;
\path{research/all-rank-nonhostile-one-cycle-theorem-2026-07-30.md}.

\bibitem{TwoPent} Companion manuscript, \emph{Positive Square Energy for Cacti
with Triangles and Two Pentagons}, 2026;
\path{all-rank-triangle-two-pentagon-cacti/paper.tex}.

\bibitem{Bouquet} Companion manuscript, \emph{Positive Square Energy for a
Two-$C_5$ Bouquet with Arbitrary Trees}, 2026;
\path{two-c5-bouquet-trees/paper.tex}.

\bibitem{Connector} Companion manuscript, \emph{Positive Square Energy for
Two Pentagons Joined by an Arbitrary Path}, 2026;
\path{two-c5-all-connectors/paper.tex}.
\end{thebibliography}

\end{document}
